\bookchapter{Capítulo 18: Backup e migração de dados}

\chapterintro{Um vizinho meu deixou o celular cair na piscina e, quando o tirou, já não ligava mais. O que mais a entristeceu não foi o celular em si, mas as fotos de três anos sem backup — todo o registro do filho, do nascimento até aprender a andar. Ouvir isso me deixou muito mal, porque era algo totalmente evitável. Neste capítulo, vamos fazer um seguro de verdade para seus dados.}

\booksection{I. Primeiro, entenda: o que no seu celular, se perder, desaparece para sempre}

O “estado de sobrevivência” de cada tipo de dado é diferente. Você precisa distinguir:

\begin{itemize}
  \item \textbf{Armazenado localmente no celular, se perder, desaparece mesmo}: Fotos e vídeos (se não tiver sincronização em nuvem ativada), notas, SMS, histórico de chamadas, arquivos baixados.
  \item \textbf{Armazenado no servidor do app, ao trocar de celular e fazer login novamente, está lá}: Lista de amigos do WeChat, pedidos do Taobao, histórico de transações no app do banco. Essas coisas estão armazenadas nos servidores das respectivas empresas, não no seu celular.
  \item \textbf{Na zona cinzenta, exige operação especial}: Histórico de conversas do WeChat (embora o servidor do WeChat armazene, você mesmo não consegue baixar do servidor), credenciais de dispositivo vinculadas ao app do banco, chaves secretas nos apps de autenticação de dois fatores (como o Google Authenticator).
\end{itemize}

Portanto, o que você precisa proteger com prioridade são a primeira e a terceira categorias.

\booksection{II. Backup completo do dispositivo: resolva tudo de uma vez}

O backup completo é como tirar um “instantâneo” completo do celular — configurações atuais do sistema, layout da tela inicial, dados de apps, fotos, contatos, tudo empacotado.

\booksubsection{Opção 1: Backup local no computador (a mais confiável, recomendada para quem tem computador)}

\begin{itemize}
  \item \textbf{iPhone}: Conecte o celular ao computador com o cabo de dados, abra o iTunes (Windows antigo) ou Finder (Mac) e, na página do dispositivo, selecione "Fazer Backup Agora". Atenção: marque a opção "Criptografar backup", caso contrário as senhas e os dados de saúde não serão incluídos no backup.
  \item \textbf{Android}: Cada marca tem seu próprio software para computador, como Huawei Phone Assistant, Xiaomi Assistant, Samsung Smart Switch. Conecte ao computador e siga as instruções.
  \item A vantagem do backup no computador é não precisar pagar por espaço em nuvem, e a restauração é muito mais rápida que pela nuvem. A desvantagem é que você precisa lembrar de conectar periodicamente para fazer o backup — e a maioria das pessoas não consegue fazer isso.
\end{itemize}

\booksubsection{Opção 2: Backup em nuvem (a mais prática, recomendada para todos)}

\begin{itemize}
  \item \textbf{iPhone}: Ajustes > Seu nome > iCloud > Backup do iCloud > Ativar. O backup automático é feito quando o celular está conectado à energia, no Wi-Fi e com a tela bloqueada. Os 5GB grátis provavelmente não serão suficientes; recomenda-se assinar 50GB.
  \item \textbf{Android}: Cada marca tem sua função de backup em nuvem. Huawei: procure por "Backup em nuvem" nas configurações; Xiaomi: "Xiaomi Cloud"; OPPO/vivo: similar. Além disso, há o backup nativo do Google (Ajustes > Sistema > Backup) — ative-o também.
  \item A maior vantagem do backup em nuvem é "não precisar se preocupar" — é automático. A maior desvantagem é que a restauração demora muito para baixar, e o espaço gratuito geralmente não é suficiente.
\end{itemize}

\booksection{III. Estratégia de backup por categoria: o que é importante, salve uma cópia extra separada}

O backup completo do dispositivo é a rede de segurança. Mas algumas coisas são tão importantes que merecem uma cópia extra separada:

\booksubsection{Fotos e vídeos (os ativos digitais mais preciosos)}

\begin{itemize}
  \item \textbf{Plano de seguro duplo}: Ative a sincronização automática do iCloud Fotos / Google Fotos e, ao mesmo tempo, a cada poucos meses, copie manualmente as fotos para o computador ou HD externo. Assim, uma cópia na nuvem e uma cópia local, uma serve de backup da outra.
  \item \textbf{Plano sem gastar dinheiro}: O Google Fotos oferece 15GB grátis; você pode configurar para "Qualidade compactada" e economizar espaço. Ou, uma vez por mês, copie todas as fotos para o computador com o cabo de dados.
\end{itemize}

\booksubsection{Contatos (fácil de ignorar, mas extremamente importante)}

\begin{itemize}
  \item Os contatos do iPhone, por padrão, sincronizam com o iCloud. Se a sincronização estiver ativada, não precisa de backup extra. Mas, para ficar ainda mais seguro, você pode acessar o iCloud.com e exportar um arquivo vCard para o computador.
  \item Os contatos do Android geralmente sincronizam com o Google ou com a nuvem da marca. Vá nas configurações do app Contatos e exporte um arquivo .vcf — isso é um backup completo.
\end{itemize}

\booksubsection{Histórico de conversas do WeChat (tratamento separado, prioridade máxima)}

Este deve ser tratado separadamente, porque o histórico do WeChat muitas vezes fica incompleto no backup completo do dispositivo, e as ferramentas de migração de celular simplesmente não conseguem transferi-lo.

\begin{itemize}
  \item \textbf{Método 1: Migrar para o novo celular}. WeChat Eu Ajustes Geral Migração e Backup de Histórico de Conversas Migrar Migrar para Celular/Tablet. Conecte o celular antigo e o novo na mesma rede Wi-Fi, escaneie o código QR e transfira. Atenção: não saia do WeChat no celular antigo até a migração terminar.
  \item \textbf{Método 2: Backup para o computador}. Instale o WeChat para PC no computador. No canto inferior esquerdo, três barras Migração e Backup Backup e Restauração Fazer Backup do Histórico de Conversas no Computador. Depois, ao trocar de celular, restaure no novo.
  \item \textbf{Método 3: Backup para a nuvem}. O WeChat também suporta backup para o iCloud ou para o espaço em nuvem do celular. Mas atenção: este backup só pode ser restaurado no mesmo celular, não permite restauração entre dispositivos. Se for trocar de celular, é preciso usar os dois primeiros métodos.
\end{itemize}

\booksection{IV. Guia de troca de celular: do antigo para o novo}

Trocar de celular é motivo de alegria, mas se os dados não forem transferidos, vira uma dor de cabeça. Trate cada caso separadamente:

\booksubsection{Upgrade na mesma marca (o mais simples)}

\begin{itemize}
  \item \textbf{iPhone Novo iPhone}: Ao ligar e ativar o novo iPhone, coloque o antigo por perto. Aparecerá automaticamente o "Início Rápido". Use o celular antigo para escanear o padrão no novo. Todos os dados, configurações, layout de apps serão transferidos exatamente como estavam, incluindo as senhas de Wi-Fi. Leva cerca de 20 a 30 minutos.
  \item \textbf{Android da mesma marca}: Use a ferramenta de migração nativa da marca. Huawei: "Phone Clone"; Xiaomi: "Mi Mover"; OPPO: "Phone Switch"; vivo: "EasyShare". Os passos são similares: instale o mesmo app nos dois celulares, conecte por QR code, selecione o que deseja transferir e inicie.
\end{itemize}

\booksubsection{Troca entre marcas diferentes (Android para Android, marcas diferentes)}

Use ferramentas de migração de terceiros ou o app de migração da marca antiga (a maioria dos apps de migração das marcas também suporta transferência para outros Androids). O que pode ser transferido: contatos, fotos, SMS, histórico de chamadas, alguns apps. O que não pode: estado de login nos apps (você precisará fazer login novamente), histórico de conversas do WeChat (use a função de migração do próprio WeChat).

\booksubsection{Troca entre plataformas (Android para iPhone, ou vice-versa)}

\begin{itemize}
  \item \textbf{Android iPhone}: A Apple tem o app "Migrar para o iOS". Baixe no Android e, ao ativar o novo iPhone, selecione "Transferir Dados do Android". Transfere contatos, fotos, SMS e outros dados básicos, mas não os dados dos apps. O histórico do WeChat precisa ser migrado pela função do próprio WeChat.
  \item \textbf{iPhone Android}: Cada marca Android tem sua ferramenta de migração correspondente, mas, sinceramente, a experiência de migração entre plataformas não é tão boa quanto na mesma plataforma. O mais importante é fazer separadamente antes: exportar contatos, fazer backup das fotos, migrar o histórico do WeChat. O resto, devagar se resolve.
\end{itemize}

\booksection{V. A regra de ouro do backup: princípio 3-2-1}

Este é um princípio clássico na área de segurança de dados, simples e fácil de lembrar:

\begin{itemize}
  \item \textbf{3}: Dados importantes devem ter pelo menos três cópias (uma original, duas de backup).
  \item \textbf{2}: Os backups devem usar dois tipos diferentes de mídia de armazenamento (por exemplo, uma na nuvem, uma em HD externo).
  \item \textbf{1}: Pelo menos uma cópia de backup deve estar em um local físico diferente (para prevenir que incêndio, roubo etc. destruam todos os dados simultaneamente).
\end{itemize}

Para uma pessoa comum, conseguir “uma cópia na nuvem e uma no computador” já é muito bom. Não precisa complicar demais. O importante é realmente fazer.

\begin{tcbwarn}{Atenção especial ao histórico de conversas do WeChat}
O histórico de conversas do WeChat não pode ser migrado por nenhuma ferramenta de migração de celular! Ele não está coberto pelo backup completo do dispositivo. É obrigatório usar a função nativa “Migração de Histórico de Conversas” do WeChat ou fazer backup para o WeChat no computador. Muita gente só descobre que perdeu todo o histórico depois de trocar de celular — aí já é tarde, pois o celular antigo pode já ter sido restaurado para as configurações de fábrica.

\end{tcbwarn}

\begin{tcbpractice}{Exercício: faça um simulado completo de backup}
\begin{enumerate}
  \item Verifique se o Backup do iCloud ou o backup em nuvem da marca está ativado e confirme a data do último backup.
  \item Faça backup manual das fotos: conecte o celular ao computador com o cabo de dados e copie as pastas DCIM e Pictures para o computador.
  \item Faça backup dos contatos: usuários iPhone exportem vCard pelo iCloud.com; usuários Android exportem arquivo .vcf nas configurações do app Contatos.
  \item Faça backup do histórico do WeChat: abra o WeChat Eu Ajustes Geral Migração e Backup de Histórico de Conversas Fazer backup no computador.
  \item Copie o arquivo vcf de contatos e a pasta de fotos exportados para um HD externo ou pen drive, completando a parte do "3-2-1": "uma cópia offline em outro lugar".
\end{enumerate}

\end{tcbpractice}
